% Generated draft: requires mathematical and source review.
\documentclass[11pt]{article}
\usepackage[margin=25mm]{geometry}
\usepackage{fontspec}
\setmainfont{Latin Modern Roman}
\usepackage{amsmath,amssymb,mathtools,unicode-math}
\setmathfont{Latin Modern Math}
\usepackage{enumitem,needspace,xurl,hyperref}
\setlength{\parindent}{0pt}
\setlength{\parskip}{5pt}
\setlength{\emergencystretch}{3em}
\newcommand{\N}{\mathbb{N}}
\newcommand{\Z}{\mathbb{Z}}
\newcommand{\Q}{\mathbb{Q}}
\newcommand{\R}{\mathbb{R}}
\newcommand{\C}{\mathbb{C}}
\newcommand{\F}{\mathbb{F}}
\newcommand{\A}{\mathbb{A}}
\newcommand{\PP}{\mathbb{P}}
\newcommand{\E}{\mathbb{E}}
\newcommand{\eps}{\varepsilon}
\newcommand{\dd}{\,\mathrm{d}}
\begin{document}
\section*{UnsolvedMath ID: 30004456}
\textbf{Source problem code:} catalogue entry
\subsection{Definitions and mathematical statement}
Let \(F_n\) be a free group of finite rank \(n\), let \(\alpha\in\operatorname{Aut}(F_n)\), and let
\[
G=F_n\rtimes_\alpha\Z.
\]
Thus, choosing a generator \(s\) of the displayed cyclic factor, one may write
\[
G=\langle F_n,s\mid sxs^{-1}=\alpha(x)\text{ for every }x\in F_n\rangle.
\]
A character here means a group homomorphism \(\phi:G\to\Z\). The supplied statement does not require \(\phi\) to be nonzero or surjective.

An HNN splitting \(G=A*_B\) means an isomorphism with a presentation
\[
\left\langle A,t\ \middle|\ t^{-1}\iota_0(b)t=\iota_1(b)\text{ for every }b\in B\right\rangle,
\]
where \(\iota_0,\iota_1:B\hookrightarrow A\) are injective homomorphisms. Its standard associated character \(\psi:G\to\Z\) satisfies \(\psi(a)=0\) for every \(a\in A\) and \(\psi(t)=1\). In particular, \(\psi\) is surjective. The source uses the phrase ``inducing \(\phi\)'' without specifying whether this means \(\phi=\psi\), equality up to sign, or a different convention involving an integral multiple of \(\psi\). This distinction is unresolved for the supplied domain of all characters.

For a group \(H\), write
\[
b_k^{(2)}(H)=\dim_{\mathcal N(H)}H_k^{(2)}(H),
\qquad
\chi^{(2)}(H)=\sum_{k\geq 0}(-1)^k b_k^{(2)}(H),
\]
where \(\mathcal N(H)\) is its group von Neumann algebra and \(H_k^{(2)}(H)\) denotes its reduced \(L^2\)-homology, whenever the alternating sum is defined under the chosen finiteness or convergence convention. Likewise, when the rational Betti numbers are finite and vanish in sufficiently high degrees, the ordinary group Euler characteristic is
\[
\chi(H)=\sum_{k\geq 0}(-1)^k\dim_{\Q}H_k(H;\Q).
\]
The supplied record specifies no convention for Euler characteristics beyond these finiteness conditions, and supplies no explicit hypotheses ensuring that every quantity in the proposed equality is defined.

\textbf{Conjectural equality as supplied.} For every free-by-cyclic group \(G=F_n\rtimes_\alpha\Z\) and every character \(\phi:G\to\Z\), is
\[
-\chi^{(2)}(\ker\phi)
=
\min\left\{
-\chi(A)\ \middle|\
\begin{gathered}
G=A*_B\text{ is an HNN splitting inducing }\phi,\\
A\text{ and }B\text{ are finitely generated}
\end{gathered}
\right\}?
\]
This is the extracted conjectural formula, with its unresolved qualifications retained. A precise assertion for the entire stated domain requires a convention for ``inducing,'' treatment of zero and non-surjective characters, conditions or conventions making both Euler characteristics meaningful, and a rule for an empty admissible set or a set whose minimum is not attained. Under the standard associated-character convention above, only surjective characters can be induced; no restriction to such characters is silently imposed here.
\subsection{Short English statement}
The proposed equality compares the negative \(L^2\)-Euler characteristic of the kernel of an integer-valued character of a free-by-cyclic group with the smallest negative Euler characteristic of a finitely generated HNN base inducing that character. The meaning of ``inducing,'' the allowed characters, and the existence of the Euler characteristics and minimum require clarification. This is a statement-only draft requiring human mathematical review.
\subsection{Sources}
\begin{itemize}
\item Local UnsolvedMath registry: \texttt{data/unsolvedmath/problems.json}, with \texttt{catalog/problems.json} as the additional registry specified by the project instructions; permanent ID \(30004456\), exact source problem code \texttt{catalogue entry}, title ``\(L^2\) Euler Characteristics and HNN Splittings.'' The supplied complete record is the evidence used for this draft; the registry files were not independently inspected.
\item Supplied report citation: \emph{Geometric Structures in Group Theory} (2021), \emph{Oberwolfach Reports}, DOI \texttt{10.4171/owr/2020/16}, pp.\ 877--918. Supplied source link: \url{https://doi.org/10.4171/owr/2020/16}. The record describes the reports as documenting mathematical workshops and their open-problem discussions. No individual proposer is supplied; the proposed year is recorded as 2020.
\item Supplied literature source: ``Oberwolfach Report 16/2020,'' \url{https://doi.org/10.4171/owr/2020/16}.
\item Literature assessment dated exactly \texttt{2026-08-22}: the supplied curation labels the problem open, states that the associated polytope and the \(L^2\)-Euler characteristic are computable, and states that the proposed equality with minimal HNN-splitting complexity is known for one-relator groups while the extracted free-by-cyclic conjecture remains open in general. These are supplied dated curation claims, not independently verified mathematical or literature conclusions. The background explicitly requests independent verification of time-sensitive status and cited claims.
\item The registry assigns the statement status \texttt{corrected\_verified}, explaining that notation, typography, omitted definitions, or duplicate clauses were repaired while preserving mathematical content. This supplied editorial label does not establish mathematical correctness or resolve the formulation issues recorded here.
\end{itemize}
\end{document}
